\documentclass[12pt,presentation]{beamer}
\usepackage[utf8]{inputenc}
\usepackage[T1]{fontenc}
\usepackage{fixltx2e}
\usepackage{graphicx}
\usepackage{longtable}
\usepackage{float}
\usepackage{wrapfig}
\usepackage{soul}
\usepackage{textcomp}
\usepackage{marvosym}
\usepackage{wasysym}
\usepackage{latexsym}
\usepackage{amssymb}
\usepackage{hyperref}
\tolerance=1000
\usepackage[english]{babel} \usepackage{ae,aecompl}
\usepackage{mathpazo,courier,euler} \usepackage[scaled=.95]{helvet}
\usepackage{listings}
\lstset{language=Python, basicstyle=\ttfamily\bfseries,
commentstyle=\color{red}\itshape, stringstyle=\color{green},
showstringspaces=false, keywordstyle=\color{blue}\bfseries}
\providecommand{\alert}[1]{\textbf{#1}}

\title{Getting started with TDD}
\author{FOSSEE}
\institute{IIT Bombay}

\usetheme{Warsaw}\usecolortheme{default}\useoutertheme{infolines}\setbeamercovered{transparent}

\AtBeginSection[]
{
  \begin{frame}<beamer>
    \frametitle{Outline}
    \tableofcontents[currentsection]
  \end{frame}
}

\begin{document}

\begin{frame}

\begin{center}
\vspace{12pt}
\textcolor{blue}{\huge Getting started with Test Driven Development}
\end{center}
\vspace{18pt}
\begin{center}
\vspace{10pt}
\includegraphics[scale=0.95]{../../images/fossee-logo.png}\\
\vspace{5pt}
\scriptsize Developed by FOSSEE Team, IIT-Bombay. \\ 
\scriptsize Funded by National Mission on Education through ICT\\
\scriptsize  MHRD,Govt. of India\\
\includegraphics[scale=0.30]{../../images/iitb-logo.png}\\
\end{center}
\end{frame}

\begin{frame}
\frametitle{Outline}
\setcounter{tocdepth}{3}
\tableofcontents
\end{frame}

\section{Introduction}

\begin{frame}
  \frametitle{Objectives}
  At the end of this section, you will be able to:
  \begin{itemize}
  \item Understand the concept of Test Driven Development.
  \item Write your code using the TDD paradigm.
  \end{itemize}
\end{frame}


\begin{frame}
\frametitle{Pre-requisite}
\label{sec-3}

Spoken tutorial on Basic Python
\begin{itemize}
\item Topic of tutorial here
\end{itemize}
\end{frame}


\begin{frame}
  \frametitle{What is TDD?}
  The basic steps of TDD are roughly as follows --
  \begin{enumerate}
  \item Decide upon the feature to implement and the methodology of
    testing it.
  \item Write the tests for the feature decided upon.
  \item Just write enough code, so that the test can be run, but it fails.
  \item Improve the code, to just pass the test and at the same time
    passing all previous tests.
  \item Run the tests to see, that all of them run successfully.
  \item Refactor the code you've just written -- optimize the algorithm,
    remove duplication, add documentation, etc.
  \item Run the tests again, to see that all the tests still pass.
  \item Go back to 1.
  \end{enumerate}
\end{frame}

\section{First Test}

\begin{frame}[fragile]
  \frametitle{First Test -- GCD}
  \begin{itemize}
  \item simple program -- GCD of two numbers
  \item What are our code units?
    \begin{itemize}
    \item Only one function \texttt{gcd}
    \item Takes two numbers as arguments
    \item Returns one number, which is their GCD
    \end{itemize}
\begin{lstlisting}
c = gcd(44, 23)
\end{lstlisting}
  \item c will contain the GCD of the two numbers.
  \end{itemize}
\end{frame}

\begin{frame}[fragile]
  \frametitle{Test Cases}
  \begin{itemize}
  \item Important to have test cases and expected outputs even before
    writing the first test!
  \item $a=48$, $b=48$, $GCD=48$
  \item $a=44$, $b=19$, $GCD=1$
  \item Tests are just a series of assertions
  \item True or False, depending on expected and actual behavior
  \end{itemize}

\end{frame}

\begin{frame}[fragile]
  \frametitle{Test Cases -- Code}
\begin{lstlisting}
tc1 = gcd(48, 64)
if tc1 != 16:
    print "Failed for a=48, b=64. Expected 16. \
    Obtained %d instead." % tc1
    exit(1)

tc2 = gcd(44, 19)
if tc2 != 1:
    print "Failed for a=44, b=19. Expected 1. \
    Obtained %d instead." % tc2
    exit(1)

print "All tests passed!"
\end{lstlisting}
\begin{itemize}
\item The function \texttt{gcd} doesn't even exist!
\end{itemize}
\end{frame}


\begin{frame}[fragile]
  \frametitle{Stubs}
  \begin{itemize}
  \item First write a very minimal definition of \texttt{gcd}
    \begin{lstlisting}
def gcd(a, b):
    pass
    \end{lstlisting}
  \item Written just, so that the tests can run
  \item Obviously, the tests are going to fail
  \end{itemize}
\end{frame}

\begin{frame}[fragile]
  \frametitle{\texttt{gcd.py}}
\begin{lstlisting}
def gcd(a, b):
    pass

if __name__ == '__main__':
    tc1 = gcd(48, 64)
    if tc1 != 16:
        print "Failed for a=48 and b=64. \
        Expected 16. Obtained %d instead." % tc1
        exit(1)
    tc2 = gcd(44, 19)
    if tc2 != 1:
        print "Failed for a=44 and b=19. \
        Expected 1. Obtained %d instead." % tc2
        exit(1)
    print "All tests passed!"
\end{lstlisting}
\end{frame}

\begin{frame}[fragile]
  \frametitle{First run}
\begin{lstlisting}
$ python gcd.py
Traceback (most recent call last):
  File "gcd.py", line 7, in <module> 
  print "Failed for a=48 and b=64. Expected 16. 
  Obtained %d instead." % tc1
TypeError: %d format: 
a number is required, not NoneType
\end{lstlisting} %$

  \begin{itemize}
  \item We have our code unit stub, and a failing test. 
  \item The next step is to write code, so that the test just passes.
  \end{itemize}
\end{frame}

\begin{frame}[fragile]
  \frametitle{Euclidean Algorithm}
  \begin{itemize}
  \item Modify the \texttt{gcd} stub function
  \item Then, run the script to see if the tests pass.
  \end{itemize}
\begin{lstlisting}
def gcd(a, b):
    if a == 0:
        return b
    while b != 0:
        if a > b:
            a = a - b
        else:
            b = b - a
    return a
\end{lstlisting}
\begin{lstlisting}
$ python gcd.py
All tests passed!
\end{lstlisting} %$
  \begin{itemize}
  \item \alert{Success!}
  \end{itemize}
\end{frame}


\begin{frame}[fragile]
  \frametitle{Euclidean Algorithm -- Modulo}
  \begin{itemize}
  \item Repeated subtraction can be replaced by a modulo
  \item modulo of \texttt{a\%b} is always less than b
  \item when \texttt{a < b}, \texttt{a\%b} equals \texttt{a}
  \item Combine these two observations, and modify the code
\begin{lstlisting}
def gcd(a, b):
    while b != 0:
        a, b = b, a % b
    return a
\end{lstlisting}
  \item Check that the tests pass again
  \end{itemize}
\end{frame}

\begin{frame}[fragile]
  \frametitle{Euclidean Algorithm -- Recursive}
  \begin{itemize}
  \item Final improvement -- make \texttt{gcd} recursive
  \item More readable and easier to understand
\begin{lstlisting}
def gcd(a, b):
    if b == 0:
        return a
    return gcd(b, a%b)
\end{lstlisting}
  \item Check that the tests pass again
  \end{itemize}
\end{frame}

\begin{frame}[fragile]
  \frametitle{Document \texttt{gcd}}
  \begin{itemize}
  \item Undocumented function is as good as unusable
  \item Let's add a docstring \& We have our first test!
  \end{itemize}
\begin{lstlisting}
def gcd(a, b):
    """Returns the Greatest Common Divisor of the 
    two integers passed as arguments.

    Args:
      a: an integer
      b: another integer

    Returns: Greatest Common Divisor of a and b
    """
    if b == 0:
        return a
    return gcd(b, a%b)
\end{lstlisting}
\end{frame}


\begin{frame}[fragile]
  \frametitle{Persistent Test Cases}
  \begin{itemize}
  \item Tests should be pre-determined and written, before the code
  \item Test Data is repeatedly used; Hence, saved in persistent
    format
  \item Let's save data for GCD tests in a text file. 
  \item The file shall have multiple lines of test data
  \item Each line corresponds to a single test case
  \item Each line consists of three comma separated values --
    \begin{itemize}
    \item First two coloumns are the integers for which the GCD has to be
      computed
    \item Third coloumn is the expected GCD to the preceding two
      numbers.
    \end{itemize}
  \item Let us call our data file \texttt{gcd\_testcases.dat}
  \end{itemize}
\end{frame}

\begin{frame}[fragile]
  \frametitle{Modify \texttt{gcd.py}}
\begin{lstlisting}
if __name__ == '__main__':
    for line in open('gcd_testcases.dat'):
        values = line.split(', ')
        a = int(values[0])
        b = int(values[1])
        g = int(values[2])

        tc = gcd(a, b)
        if tc != g:
            print "Failed for a=%d and b=%d.\
            Expected %d. Obtained %d instead."\
            % (a, b, g, tc)
            exit(1)

    print "All tests passed!"
\end{lstlisting}
\end{frame}

\begin{frame}

  \begin{block}{}
  \begin{center}
  \textcolor{blue}{\Large THANK YOU!} 
  \end{center}
  \end{block}
\begin{block}{}
  \begin{center}
    For more Information, visit our website\\
    \url{http://fossee.in/}
  \end{center}  
  \end{block}
\end{frame}

\end{document}
